\documentclass[12pt]{article}

\usepackage[utf8]{inputenc} % important for pdflatex + any UTF-8 characters
\usepackage{amsmath, amssymb}
\usepackage{geometry}
\geometry{margin=1in}

\setlength{\parindent}{0pt}

\title{Practice Problem Set 1}
\author{ECON 300 -- Labour Economics}
\date{Winter 2026}

\begin{document}
\maketitle



\subsection*{Problem 1: Labour Force Participation}

In a city with a working-age population (15+) of 800,000:

\begin{itemize}
    \item Employed: 520,000
    \item Unemployed: 40,000
    \item Not in the labour force: 240,000
\end{itemize}

\begin{enumerate}
    \item Calculate the \textbf{labour force}.
    \item Compute the \textbf{labour force participation rate (LFPR)}.
    \item Compute the \textbf{unemployment rate}.
    \item Compute the \textbf{employment-to-population ratio}.
    \item Interpret the difference between the LFPR and the employment-to-population ratio.
\end{enumerate}

\bigskip\bigskip

\subsection*{Problem 2: Labour Supply and the Reservation Wage}

A worker has the following weekly budget constraint:
\[
C = 500 + 20h
\]
where \(C\) is consumption (dollars) and \(h\) is hours worked. The worker has \$500 in non-labour income and earns \$20 per hour.

The worker's utility function is:
\[
U = C - 0.5h^2
\]

\begin{enumerate}
    \item Derive the worker's \textbf{labour supply function}.
    \item Find the optimal number of hours worked.
    \item Determine the \textbf{reservation wage}.
    \item Explain the economic meaning of the reservation wage.
\end{enumerate}

\bigskip\bigskip

\subsection*{Problem 3: Income and Substitution Effects}

Suppose a worker earns \$25 per hour and works 40 hours per week. The wage increases to \$30 per hour.

\begin{enumerate}
    \item Calculate the change in weekly income.
    \item Using the labour--leisure model, explain:
    \begin{itemize}
        \item the \textbf{substitution effect}, and
        \item the \textbf{income effect}.
    \end{itemize}
    \item Under what conditions could hours worked fall after a wage increase?
\end{enumerate}

\bigskip\bigskip

\subsection*{Problem 4: Labour Supply and Leisure (Graphical Analysis)}

Draw a standard \textbf{labour--leisure diagram} with leisure on the horizontal axis and consumption on the vertical axis.

\begin{enumerate}
    \item Show the budget constraint for a wage of \$20 per hour.
    \item Illustrate an increase in the wage to \$30 per hour.
    \item Decompose the change in hours worked into:
    \begin{itemize}
        \item substitution effect, and
        \item income effect.
    \end{itemize}
    \item Label the \textbf{reservation wage} on the graph.
\end{enumerate}

\bigskip\bigskip

\subsection*{Problem 5: Market Labour Supply and Demand}

Assume the labour market is described by:
\[
L^D = 100 - 2w
\qquad
L^S = 20 + 3w
\]
where \(w\) is the wage.

\begin{enumerate}
    \item Find the \textbf{equilibrium wage} and \textbf{employment}.
    \item Graph the labour supply and demand curves.
    \item Suppose a minimum wage of \$18 is imposed:
    \begin{itemize}
        \item Is it binding?
        \item Calculate unemployment.
    \end{itemize}
    \item Explain the result using economic intuition.
\end{enumerate}

\end{document}
